\subsection{有限域上的单代数}

定理 2（韦德伯恩）. —— 每个有限体都是交换的。

设 D 是一个有限体，K 是它的中心。K-代数 D 是一个次数为 \(m^{2}\) 的中心单代数，其中 m 是一个严格正整数。约化范数是一个次数为 m 的齐次多项式映射 \(Nrd : D \to K\)（VIII, 第 345 页，命题 6），且我们有 \(\text{Nrd}(a) \neq 0\)（对 D 中每个 \(a \neq 0\)；VIII, 第 340 页，命题 3）。上面的推论蕴含 \(m \geqslant m^{2}\)，因而 m = 1。于是我们有 D = K。

推论 1. —— 每个有限单环都同构于一个矩阵环 \(\mathbf{M}_{n}(\mathrm{L})\)，其中 n 是一个严格正整数，L 是一个有限交换域。

这由定理 2 与单环的结构定理（VIII, 第 120 页，定理 1）得出。

推论 2. —— 设 K 是一个有限交换域。K 上每个中心单代数都同构于一个矩阵代数 \(\mathbf{M}_{n}(\mathrm{K})\)，其中 n 是一个严格正整数。

这由定理 2 与中心单代数的结构定理（VIII, 第 252 页，定理 1）得出。

注. —— 1) 这里给出定理 2 的另一个证明。设 D 是一个有限体，K 是它的中心，L 是 D 的一个极大交换子域。设 x 是 \(D^{*}\) 的一个元素。它属于 D 的一个极大交换子域 \(L_{1}\)。我们有等式

\[
[ \mathrm{D}: \mathrm{K} ] = [ \mathrm{L}: \mathrm{K} ] ^ {2} = [ \mathrm{L} _ {1}: \mathrm{K} ] ^ {2}
\]

这由 VIII, 第 265 页，推论 2 得出，因而 \([\mathrm{L}:\mathrm{K}] = [\mathrm{L}_1:\mathrm{K}]\)。由 V, §12, No.~2, 第 94 页，命题 3，K 的扩张 L 与 \(\mathrm{L}_1\) 是同构的。由 VIII, 第 263 页，推论，存在一个元素 \(a\)（属于 \(\mathrm{D}^*\)）使得 \(a\mathrm{La}^{-1} = \mathrm{L}_1\)，于是 \(a^{-1}xa\) 属于 L。于是我们有 \((ay)^{-1}x(ay) = a^{-1}xa\)（对每个 \(y\in \mathrm{L}^*\)）。因此，若 S 是 \(\mathrm{D}^*\) 模 \(\mathrm{L}^*\) 的左陪集的一个代表系，则 \(\mathrm{D}^* -\{1\}\) 的每个元素都可写成 \(sxs^{-1}\)（其中 \(s\in \mathrm{S}\)，\(x\in \mathrm{L}^* -\{1\}\)）。我们把 \(\mathrm{D}^*\) 的阶记作 \(d\)，把 \(\mathrm{L}^*\) 的阶记作 \(\ell\)。由于 S 的基数等于 \(d / \ell\)，我们有 \(d - 1\leqslant (d / \ell)(\ell -1) = d - d / \ell\)。由此推出 \(\ell = d\)，因而 \(\mathrm{L} = \mathrm{D}\)，这就证明体 D 是交换的。

\begin{enumerate}
\def\labelenumi{\arabic{enumi})}
\setcounter{enumi}{1}
\tightlist
\item
  设 L 是一个具有下列性质的交换域：
\end{enumerate}

(C \(_{1}\) ) 设 V 是域 L 上一个维数为 n 的向量空间，d 满足 0 \textless{} d \textless{} n。~对每个次数为 d 的齐次多项式映射 F : V → L，存在 V 的一个非零元素 x 使得 F(x) = 0。

定理 2 的证明表明，每个以 L 为中心、在 L 上有限次数的体都等于 L。

由 VIII, 第 356 页的推论，每个有限域都具有性质 \((\mathrm{C}_{1})\)。我们可以证明（VIII, 第 360 页，习题 7）下列域具有性质 \((\mathrm{C}_{1})\)：

- 有限域的每个代数扩张

-- 代数闭域上单变量有理分式所成的域（曾定理）。

- *每个赋有离散赋值、对该赋值完备且其剩余域为代数闭域的域（VIII, 第 332 页，习题 17）。*

\begin{enumerate}
\def\labelenumi{\arabic{enumi})}
\setcounter{enumi}{2}
\tightlist
\item
  设体 \(\mathbf{K}\) 满足下列条件：
\end{enumerate}

- 若 \(\mathrm{L}\) 是 \(\mathrm{K}\) 的一个有限次数扩张，则它是循环的，且范数映射 \(\mathrm{N}:\mathrm{L}^{*}\to \mathrm{K}^{*}\) 是满射。

特别地，当域 K 有限时该条件满足（V, §12, No.~2, 第 95 页，命题 4）。这时我们可以证明，每个以 K 为中心、在 K 上有限次数的体都等于 K（VIII, 第 329 页，习题 10）。

